\documentclass[12pt,a4paper,oneside]{article}
\usepackage[brazil]{babel}
\usepackage[utf8]{inputenc}
\usepackage[dvipsnames]{xcolor}
\usepackage{listings}
\usepackage{wrapfig}
\usepackage{olymp}
\setcounter{problem}{1}

\begin{document}

\begin{problem}{Área do triângulo}{area}{2}

Um triângulo é uma figura geométrica formado por três lados. Pode ser definido também como a união de três pontos não-colineares por três segmentos de reta. Por exemplo, na figura da esquerda temos um triângulo de lados com tamanhos $3$, $4$ e $5$, e na da direita temos um triângulo formado pelos pontos $P = (2,1) $, $Q = (1,4)$ e $R = (6,2)$, que são unidos pelos segmentos $a$, $b$ e $c$. 

\includegraphics[scale=0.8]{images/exercise_71_0.png}\hspace{2cm}
\includegraphics[scale=0.8]{images/exercise_71_1.png}

A área do triângulo pode ser calculada pela fórmula $área = (base \times altura) / 2$. A área do triângulo da esquerda é então ($3\times4)/2=6$. Usar essa fórmula no triângulo da direita é mais complicado, porque teríamos que definir uma base e calcular a altura antes de usar a fórmula.

Outra fórmula interessante usa apenas o tamanho dos segmentos que formam o triângulo:

$área = \sqrt{s(s-a)(s-b)(s-c)}$,
sendo $s$ o valor do semiperímetro: $s = (a+b+c)/2$.

Para o triângulo da esquerda temos $s = (3+4+5)/2 = 6$ e área $\sqrt{6(6-3)(6-4)(6-5)}=6$. No da direita temos $a = \sqrt{10}$, $b = \sqrt{29}$ e $c=\sqrt{17}$, que resulta num semiperímetro $s \approx 6,33527$ e área $6.5$. Faça um programa que usa esta fórmula para calcular a área de um triângulo.

%\Notes
%
%Uma seleção ganha 3 pontos para cada vitória e 1 ponto para cada empate. Derrotas não alteram sua pontuação.

% \begin{center}
% \texttt{time1 placar1 x placar2 time2}
% \end{center}

% Onde:
% \begin{itemize}
%     \item \texttt{time1} e \texttt{time2} são os nomes das seleções (sequências de caracteres sem espaços);
%     \item \texttt{placar1} é o número de gols feitos pelo \texttt{time1};
%     \item \texttt{placar2} é o número de gols feitos pelo \texttt{time2};
% \end{itemize}


% \Notes

Use a seguinte estrutura para armazenar as coordenadas de um ponto:

\lstset{language=C++,
                basicstyle=\ttfamily,
                keywordstyle=\color{Mulberry}\ttfamily,
                stringstyle=\color{Maroon}\ttfamily,
                commentstyle=\color{YellowGreen}\ttfamily,
                morecomment=[l][\color{magenta}]{\#},
                basewidth = {.48em}
}
\begin{lstlisting}
struct Ponto {
    int x,y;
};

\end{lstlisting}

e uma função com o cabeçalho abaixo para calcular a área:

\lstset{language=C++,
                basicstyle=\ttfamily,
                keywordstyle=\color{Mulberry}\ttfamily,
                stringstyle=\color{Maroon}\ttfamily,
                commentstyle=\color{YellowGreen}\ttfamily,
                morecomment=[l][\color{magenta}]{\#},
                basewidth = {.48em}
}
\begin{lstlisting}
double area (Ponto P, Ponto Q, Ponto R);
\end{lstlisting}

Você pode também criar uma função calcular o tamanho dos seguimentos:

\lstset{language=C++,
                basicstyle=\ttfamily,
                keywordstyle=\color{Mulberry}\ttfamily,
                stringstyle=\color{Maroon}\ttfamily,
                commentstyle=\color{YellowGreen}\ttfamily,
                morecomment=[l][\color{magenta}]{\#},
                basewidth = {.48em}
}
\begin{lstlisting}
double distancia (Ponto A, Ponto B);
\end{lstlisting}

\InputFile

A entrada contém três linhas, cada uma delas com dois valores inteiros, representando as coordenadas cartesianas dos três pontos. Restrições: os valores são inteiros entre $-100$ e $100$.

\OutputFile

Seu programa deve gerar uma linha na saída contendo a área do triângulo formado pelos três pontos da entrada com 2 casas decimais.

% \Notes
% Preste atenção aos tipos das variáveis, pois $F(90) \approx 2^{61}$.

\Example


\begin{example}
\exmp{
0 0
3 0
0 4
}{
6.00
}%
\end{example}


\begin{example}
\exmp{
2 1
1 4
6 2
}{
6.50
}%
\end{example}


\begin{example}
\exmp{
-6 0
0 0
0 8
}{
24.00
}%
\end{example}


\begin{example}
\exmp{
-10 11
21 37
33 -54
}{
1566.50
}%
\end{example}

\textbf{Obs.}: os dois primeiros exemplos são os triângulos mostrados nas figuras.

%Comentários sobre o exemplo, se necessário.

\end{problem}

\end{document}
